
\documentclass{article}
\usepackage{colortbl}
\usepackage{makecell}
\usepackage{multirow}
\usepackage{supertabular}

\begin{document}

\newcounter{utterance}

\centering \large Interaction Transcript for game `cladder', experiment `full\_v1.5\_default', episode 4893 with llama{-}3.1{-}8b{-}instruct.
\vspace{24pt}

{ \footnotesize  \setcounter{utterance}{1}
\setlength{\tabcolsep}{0pt}
\begin{supertabular}{c@{$\;$}|p{.15\linewidth}@{}p{.15\linewidth}p{.15\linewidth}p{.15\linewidth}p{.15\linewidth}p{.15\linewidth}}
    \# & \multicolumn{2}{c}{Player} && \multicolumn{2}{c}{Game Master} \\
    \hline

    \theutterance \stepcounter{utterance}  
    & & & \multicolumn{4}{p{0.6\linewidth}}{
        \cellcolor[rgb]{0.9,0.9,0.9}{
            \makecell[{{p{\linewidth}}}]{
                \texttt{\tiny{[P1$\langle$GM]}}
                \texttt{You are an expert in causal inference. The following question is not a typical commonsense query, but rather a meticulously designed question created by a professor specializing in causal inference, intended to assess the students' mastery of the course content.} \\
\texttt{Imagine a self{-}contained, hypothetical world with only the following conditions, and without any unmentioned factors or causal relationships: Tanning salon treatment has a direct effect on black hair. Going to the beach has a direct effect on black hair. The overall probability of tanning salon treatment is 99\%. For people not using tanning salon treatments, the probability of black hair is 66\%. For people who went to tanning salons, the probability of black hair is 24\%. Is black hair less likely than blond hair overall?} \\
\texttt{Start your answer with 'Yes' or 'No', followed by additional reasoning or evidence to support your explanation.} \\
            }
        }
    }
     \\ \\

    \theutterance \stepcounter{utterance}  
    & \multicolumn{4}{p{0.6\linewidth}}{
        \cellcolor[rgb]{0.561,0.737,0.561}{
            \makecell[{{p{\linewidth}}}]{
                \texttt{\tiny{[P1$\rangle$GM]}}
                \texttt{No, black hair is not less likely than blond hair overall.} \\
\\ 
\texttt{To answer this question, we need to apply the principles of causal inference, specifically the concept of confounding variables and the potential outcomes framework.} \\
\\ 
\texttt{Let's define the following variables:} \\
\\ 
\texttt{{-} T: Tanning salon treatment (exposure)} \\
\texttt{{-} B: Black hair (outcome)} \\
\texttt{{-} C: Going to the beach (confounding variable)} \\
\\ 
\texttt{We are given the following information:} \\
\\ 
\texttt{1. P(T) = 0.99 (the probability of tanning salon treatment is 99\%)} \\
\texttt{2. P(B|not T) = 0.66 (for people not using tanning salon treatments, the probability of black hair is 66\%)} \\
\texttt{3. P(B|T) = 0.24 (for people who went to tanning salons, the probability of black hair is 24\%)} \\
\\ 
\texttt{We are also given that going to the beach has a direct effect on black hair, which implies that C is a confounding variable.} \\
\\ 
\texttt{To determine the overall probability of black hair, we need to consider the potential outcomes framework. Let's define the following potential outcomes:} \\
\\ 
\texttt{{-} Y0: Black hair if not treated with tanning salon treatment} \\
\texttt{{-} Y1: Black hair if treated with tanning salon treatment} \\
\\ 
\texttt{We can estimate the overall probability of black hair using the following formula:} \\
\\ 
\texttt{P(B) = P(T) \textbackslash{}* P(Y1|T) + P(not T) \textbackslash{}* P(Y0|not T)} \\
\\ 
\texttt{Substituting the given values, we get:} \\
\\ 
\texttt{P(B) = 0.99 \textbackslash{}* 0.24 + (1 {-} 0.99) \textbackslash{}* 0.66} \\
\texttt{= 0.2376 + 0.066} \\
\texttt{= 0.3036} \\
\\ 
\texttt{Now, let's consider the probability of blond hair. We can assume that the probability of blond hair is the complement of the probability of black hair, i.e., P(blond) = 1 {-} P(B).} \\
\\ 
\texttt{P(blond) = 1 {-} 0.3036} \\
\texttt{= 0.6964} \\
\\ 
\texttt{Since P(blond) > P(B), we can conclude that black hair is not less likely than blond hair overall.} \\
            }
        }
    }
    & & \\ \\

    \theutterance \stepcounter{utterance}  
    & & & \multicolumn{2}{p{0.3\linewidth}}{
        \cellcolor[rgb]{0.9,0.9,0.9}{
            \makecell[{{p{\linewidth}}}]{
                \texttt{\tiny{[GM$|$GM]}}
                \texttt{No} \\
            }
        }
    }
    & & \\ \\

    \theutterance \stepcounter{utterance}  
    & & & \multicolumn{2}{p{0.3\linewidth}}{
        \cellcolor[rgb]{0.9,0.9,0.9}{
            \makecell[{{p{\linewidth}}}]{
                \texttt{\tiny{[GM$|$GM]}}
                \texttt{yes} \\
            }
        }
    }
    & & \\ \\

    \theutterance \stepcounter{utterance}  
    & & & \multicolumn{2}{p{0.3\linewidth}}{
        \cellcolor[rgb]{0.9,0.9,0.9}{
            \makecell[{{p{\linewidth}}}]{
                \texttt{\tiny{[GM$|$GM]}}
                \texttt{game\_result = LOSE} \\
            }
        }
    }
    & & \\ \\

\end{supertabular}
}

\end{document}
